%%%%%%%%%%  scan idx 55  |  folio 42  %%%%%%%%%%
\textbf{Corollaire.} \emph{Soit $p=\dim X-\dim Y$. Alors, si $i<p$, l'homomorphisme $A^i(X)\longrightarrow A^i(U)$ est bijectif. Tout élément du noyau de $A^p(X)\longrightarrow A^p(U)$ est combinaisons linéaire de composantes de dimension $n-p$ de $Y$ (donc est proportionnel à $\ellc(Y)$ si $Y$ est irréductible)}

\subsection*{\S 2. Définition des classes de Chern des faisceaux algébriques cohérents}

Soit $\cC$ une catégorie abélienne, et soit $\cC_0$ une sous-classe de $\cC$. On désigne par $K(\cC_0)$ le groupe quotient du groupe libre engendré par les classes (à isomorphisme près) d'éléments de $\cC_0$, par le sous-groupe engendré par les éléments de la forme $(E)-(E')-(E'')$, où $E$ est une \emph{extension} de $E'$ par $E''$ ($E,E',E''\in\cC_0$). On doit supposer que les classes d'éléments de $\cC_0$ pour la relation d'isomorphisme forment un ensemble. On désigne par $\gamma_{\cC_0}(E)$ la classe d'un $E\in\cC_0$ dans $K(\cC_0)$ et l'on omet l'indicé $\cC_0$ quand il n'y a pas de risque de confusion. Le groupe $\underline{K}(\cC_0)$ est solution d'un \emph{problème universel} relatif aux fonctions $E\longrightarrow f(E)$ définies sur $\cC_0$, à valeurs dans un groupe abélien arbitraire $A$, et qui sont \emph{additives} en ce sens que $f(E)=f(E')+f(E'')$ chaque fois que $E$ est extension de $E'$ par $E''$.

Soient $F$ un foncteur additif de $\cC$ dans $\cC'$ et $\cC'_0$ une sous-classe de $\cC'$ satisfaisant à la même condition que $\cC_0$, tels que pour tout $A\in\cC_0$, les $R^qF(A)$ soient dans $\cC'_0$ et nuls pour $q$ assez grand (on suppose que dans $\cC$ existent des résolutions injectives, de sorte que les foncteurs dérivés $R^qF$ de $F$ existent)~; alors $F$ définit un homomorphisme de $\underline{K}(\cC_0)$ dans $\underline{K}(\cC'_0)$, noté encore $F$, par la formule
\begin{equation}\tag{2.11}
 F(\gamma_{\cC_0}(E))=\sum_{i\geq 0}(-1)^i\gamma_{\cC'_0}(R^iF(E)).
\end{equation}
Plus généralement, si on a un foncteur cohomologique $(F^i)$ de $\cC$ dans $\cC'$ 
